\BourbakiBackSection{历史注记}

（括号中的数字指本注记末尾的参考文献。）

对量的每一次测量都蕴含着实数的一种模糊观念（其确切理由我们将在第五章 § 2 中看到）。从数学的观点看，实数理论的起源可以追溯到巴比伦人逐步形成的一套记数法，它（原则上）能够以任意逼近程度表示任何实数{[}1{]}。掌握了这样一套记数法，以及由此自然产生的对数值计算的信心，不可避免地导致了一种“朴素”的实数观念；它与今天在初等教育中以及物理学家和工程师中间通行的、（与十进记数法相联系的）那种观念几乎毫无差别。这一观念无法精确定义，但可以这样来表述：一个数被看作是由找到它的近似值并在计算中使用这些近似值的可能性所定义的；这必然意味着，在物理量的测量（它们当然不可能承受无穷序列越来越接近的近似）与诸如 \(\sqrt{2}\) 这样的“数”（假定掌握了一种算法，能对这些数作出无穷序列越来越接近的近似）之间存在着某种混淆。

类似的“实用主义”态度出现于一切计算技巧重于严谨性与理论的数学学派之中。然而在希腊数学中占主导地位的恰恰是后者；正是希腊人给我们留下了第一个严格而连贯的量的比的理论，它本质上就是实数的理论。这一理论是关于比例、特别是不可公度比的一系列发现的顶峰；这些发现在希腊思想史上的重要性无论怎样强调都不为过，但由于缺乏准确的文献，只能看出其大致轮廓。早期阶段的希腊数学与关于比例、相似与比的种种思索密不可分地纠缠在一起，这些思索一部分是科学的，一部分是哲学的与神秘主义的，尤其是关于“简单比”（可用分子分母都很小的分数表示的比）；而毕达哥拉斯学派的特征倾向之一，便是试图用整数与整数之比来解释一切。然而，恰恰是毕达哥拉斯学派发现了正方形的对角线与其边不可公度（换句话说，\(\sqrt{2}\) 是无理数）。这无疑是数学中不可能性证明的第一个例子，而仅仅提出这样一个问题这件事本身，就意味着比与它的近似值之间被明确地区分开来，并显示出把希腊数学家与其前辈分隔开来的巨大鸿沟（*）。

关于伴随并接续这一重要发现的思想运动，我们知之甚少（ \(^{**}\) ）。我们只打算简要概述量的比理论赖以建立的主要思想；这一理论由伟大的数学家 Eudoxus（Plato 的同时代人与朋友）构造，被古典时期的希腊数学最终采纳，并经由 Euclid 的《几何原本》{[}2{]} 为我们所知，在那里（《几何原本》第五卷中）它得到了精妙的阐述：

\begin{enumerate}
\def\labelenumi{\arabic{enumi})}
\item
  “数”这个词和数的观念严格地保留给自然整数 \textgreater1 （1 是单子，严格说来并不是数），这不仅排除了我们的无理数，而且排除了我们所说的有理数：对古典时期的希腊数学家而言，后者是数之比。这里远不止一个单纯的术语问题：“数”这个词对希腊人（直到不久之前对近代人也一样）是与具有两个合成法则（加法与乘法）的体系这一观念联系在一起的；古典希腊数学家把整数之比看作算子，它定义在整数集或其某个子集上{[}p 与 q 之比是这样的算子：当 N 是 q 的倍数时，它作用于 N 便给出整数 \(p.(N/q)\) {]}，这些算子构成一个乘法群，但不构成具有两个合成法则的体系。在这一点上，希腊数学家自觉地与“计算师”即职业计算者划清界限，后者像其埃及与巴比伦前辈一样，毫无顾忌地把分数当作数来处理，或把分数加到整数上。而且，对数的概念的这一自我设限，看来与其说出于数学动机，不如说出于哲学动机，它源自最初希腊思想家关于单位与倍数的思考；单位（在这一思想体系中）不能细分，否则便失去其作为单位的特性（*）。
\item
  量的理论以公理为基础，这些公理同时适用于一切类型的量（有迹象表明，更早的理论显然是把长度、面积、体积、时间等分别处理的）。同类型的量由下述事实刻画：它们可以比较（也就是说，假设定义了相等关系——它是一个等价关系——以及关系 \textgreater{} 与 \textless{} ），它们可以加、减（A + B 有定义，且当 A \textgreater{} B 时 A - B 也有定义），并且它们满足“Archimedes 公理”（§ 2 的定理 1）。从一开始人们就清楚地认识到，最后这一事实是整座大厦的拱顶石（事实上，在实数的任何公理化刻画中它都是不可缺少的；参看第五章 § 2）。把它归属于 Archimedes 纯属偶然：在其《抛物线求积》{[}3{]} 的引言中，Archimedes 强调，这条公理曾被他的前辈们使用过，它在 Eudoxus 的工作中起过本质作用，而且由它推出的结论，其确定性与不借助它而进行的面积和体积的确定相比毫不逊色（**）。
\end{enumerate}

我们将在第五章 § 2 中看到，实数理论如何是这一公理基础的必然推论。请注意，对 Eudoxus 而言，给定类型的量构成一个具有单一内合成法则（加法）的体系，但这个体系还有一个外合成法则，其算子是量的比，这些比被设想为构成一个交换乘法群。若 A 与 A' 是同类型的量，且 B 与 B' 是（同类型的）量，则当对一切整数 m 与 m'，关系 mA \textless{} m'A' 蕴涵 mB \textless{} m'B' 且 mA \textgreater{} m'A' 蕴涵 mB \textgreater{} m'B' 时，A 与 A' 之比及 B 与 B' 之比定义为相等。比之间的不等式用类似的方法定义。这些比针对每一种类型的量都构成一个算子域，这一事实等价于第四比例项存在的公理（Euclid 的阐述中未明确陈述，但经常使用）：若给定了比 A/A'，又给定了 B，则存在与 B 同类型的 B'，使 B/B' = A/A'。这样，Eudoxus 的卓越思想使他得以把由不同类型的量所定义的一切算子域等同起来（*）；类似地，整数之比的集合（见上文）可以等同于量的比所成之集的一个子集，即有理比（或可公度的量之比）所成之集；尽管如此，由于这些比作为整数的算子（一般）只在整数集的一个子集上有定义，仍有必要单独发展它们的理论（Euclid 第七卷）。

这样构造出来的普遍算子域，对希腊数学家而言，相当于实数集之于我们；而且，凭借量的加法与量的比的乘法，他们显然已拥有相当于实数域之于我们的东西，尽管其形式远不便于处理（ \(^{**}\) ）。另一方面，人们可以问，他们是否把这些集合（给定类型的量的集合，或量的比所成之集）看作我们意义下的完备集合。否则就难以理解，他们为什么会假定第四比例项的存在（甚至没有察觉有必要把它列为公理）；此外，某些文献似乎已经涉及这类观念；他们也确实把如下事实视为自明而加以假定：可用连续运动描绘的一条曲线，要从直线的一侧过渡到另一侧，就不可能不与该直线相交。这一原理曾用于例如关于倍立方问题的研究（通过曲线的交作出 \(\sqrt[3]{2}\) ），而且它本质上等价于上述性质；尽管如此，流传至今的文献并不能使我们完全准确地辨认他们在这个问题上的想法。

这就是希腊数学古典时期实数理论的状况。Eudoxus 的构造尽管令人赞叹，在严谨性与连贯性方面无可指摘，然而必须承认，它缺乏灵活性，不利于数值计算的发展，更不利于代数计算的发展。况且，它的逻辑必然性只有对热爱严谨并熟习抽象的人才是明显的；因此，随着希腊数学的衰落，凭借计算师们的传统而保存下来的“朴素”观点逐渐重新抬头，是很自然的。这一观点例如在 Diophantus {[}4{]} 中占主导地位；他实际上与其说是官方希腊科学的拥护者，不如说是这一传统的拥护者；他敷衍了事地复述了 Euclid 的数的定义，但实际上他用“数”这个词表示代数问题中的未知数，其解可以是整数、分数或无理数（*）。尽管关于数的态度上的这一转变与数学史上最重要的进展之一即代数的发展相联系，它本身当然并不构成一种进步，而毋宁说是一种退步。

我们无法在此追溯数的概念在印度、阿拉伯与西方数学中的变迁，直到中世纪结束为止。“朴素”的数的观念占据了统治地位，而且尽管 Euclid 的《几何原本》在这一时期是数学教学的基础，Eudoxus 的学说多半一直未被理解，因为对它的需要已不再被人领会。Euclid 的“比”习惯上被说成“数”，整数的计算规则被应用到它们身上，而对这种方法何以成功却从未有人试图分析。

然而我们看到，R. Bombelli 早在十六世纪中叶就在其《代数学》{[}5{]}（*）中就这一问题阐述了一种本质上正确的观点（前提是假定 Euclid《几何原本》第五卷的结果为已知）；他认识到，一旦选定了长度单位，长度与量的比之间就有一个一一对应，于是他定义了长度之间的各种代数运算（当然假定单位已经固定），并以长度表示数，从而得到实数域的几何定义（这一观点通常被归功于 Descartes），这样就给他的代数学提供了一个坚实的几何基础（**）。

但是，Bombelli 的《代数学》尽管相对于其时代先进得出奇，却没有超出开根式以及二、三、四次方程的根式求解；当然，开根式的可能性是被不加讨论地承认的。Simon Stevin {[}6{]} 在数的问题上采取类似的立场；对他而言，数表示量的度量，并且被看作本质上是“连续的”（他没有赋予这个词以确切的含义）。他之所以区分“几何数”与“算术数”，只是出于其定义方式的偶然情况，而不是因为本性上的差别。他关于这一问题最后的话如下：“因此我们断言，根本不存在什么荒谬的、无理的、不规则的、无法解释的或不尽根的数，相反，它们之中蕴藏着如此的卓越与和谐，使我们对它们令人赞叹的完美日夜冥想不已”（{[}6{]}，第 10 页）。另一方面，他第一个把十进分数用作一种计算方法，并为它们提出了一种与我们的记法十分接近的记法；他清楚地看到，这些分数提供了一种可以无限接近任何实数的算法，其 1594 年的《代数附录》——“包含一个适用于一切方程的法则”——便表明了这一点（这本小册子已知的唯一抄本于 1914 年在 Louvain 被焚毁；但可参看{[}6{]}，第 1 卷，第 88 页）。把这样一个方程写成 \(\mathrm{P}(x) = \mathrm{Q}(x)\) 的形式{[}其中多项式 \(\mathbf{P}\) 的次数大于多项式 \(\mathbf{Q}\) 的次数，且 \(\mathrm{P(0)} < \mathrm{Q(0)}\) {]}之后，他对 \(x\) 依次代入 10、100、1000、\ldots{} ，直到发现 \(\mathrm{P}(x) > \mathrm{Q}(x)\) ，他说，这样就定出了根的位数；然后（例如设根有两位）他代入 10、20、\ldots{} ，这便定出十位数字；随后同样地定出个位数字，再依次定出各十进数字：“照这样无限地进行下去”，他说，“我们就无限地接近所要求的数”（{[}6{]}，第 88 页）。正如我们所看到的，Stevin 恰恰具有 § 6 定理 2 的思想（而且无疑是第一个具有这一思想的人），并且认识到这一定理是系统地求解数值方程的根本工具；同时我们可以看出，他对数值连续统的直观概念是如此清晰，以致要把它最终精确化已经所剩无几。

然而，在其后的两个世纪中，正确方法的最终确立两次被两种理论的发展所延误，它们的历史不属于本注记的范围：无穷小演算与级数理论。在由它们引起的种种讨论中，我们像在数学史的一切时期中一样，察觉到两类人之间的持久平衡：一类人力图冒些不甚牢靠的风险向前推进，相信以后总来得及巩固已占领的阵地；另一类人具有批判的头脑（他们在直观天赋与发明才能方面未必有丝毫逊色），相信把自己的精力用于精确表述并严格论证自己的构想并非虚掷。十七世纪争论的主要话题是“无限”小的观念，它虽然借助其所得结果而获得事后的辩护，却似乎与 Archimedes 公理公开对立；我们看到这个时期最开明的头脑最终采取了一种与 Bombelli 相差无几的观点，其区别主要在于对古人的严格方法给予了更多的注意。Isaac Barrow（Newton 的老师，他本人在无穷小演算的创立中起过重要作用）在其 1664-5-6 年间于 Cambridge 所作的数学讲演{[}7{]}中对这一观点作了出色的阐述；他认识到必须回到 Eudoxus 的理论，以便在数的问题上重新获得世所公认的“几何确定性”，他对 Eudoxus 的理论（据 Barrow 说，它对他的许多同时代人来说一直无法理解）作了长篇而极为得当的辩护，以反驳那些指责它晦涩乃至荒谬的人。另一方面，把数定义为表示量的比的符号、且可作算术运算的组合，Barrow 便以一些措辞得到了实数域，Newton 在其《算术》中重新采用了这些措辞，而直到 Dedekind 与 Cantor，他的后继者们都没有改动它们。

但正是在这个时期，级数展开的方法被引入；在死心塌地的代数学家手里，它很快带上了一种纯形式的特征，并把数学家的注意力从收敛性问题上引开，而要在实数领域内健全地使用级数，这些问题是本质性的。Newton 是这一方法的主要创立者，他当然意识到考虑这些问题的必要性；即便他没有充分阐明它们，他至少知道，他所引入的幂级数对变量的小值“通常”至少像几何级数（其收敛性古已知之）那样收敛{[}8{]}。大约同一时期，Leibniz 已注意到，各项绝对值递减且趋于 0 的交错级数是收敛的。在其后一个世纪，d'Alembert 于 1768 年对不收敛级数的使用提出了怀疑；但 Bernoulli 家族尤其是 Euler 的权威，使这类怀疑在这一时期只是罕见的例外。

显然，习惯于把级数用于数值计算的数学家绝不会如此忽视收敛概念；而在这一领域如同在许多其他领域一样，率先引导回归正确方法的是一位早年醉心于数值计算的数学家，这并非偶然。C.F. Gauss 几乎还在孩提时代就实践过算术几何平均的算法（*），他不可能不形成清晰的极限概念；在一份注明 1800 年（但直到我们这个时代才首次发表）的残稿中（{[}9{]}，第 X¹ 卷，第 390 页），他给出了实数序列的上确界与下确界以及上极限与下极限的精确定义；至于（有界序列的）确界的存在性，他似乎是视作显然而予以承认的；而上极限与下极限，他正确地定义为当 n 趋于 +∞ 时 sup \(u_{n+p}\) 与 inf \(u_{n+p}\) 的极限。另一方面，在其 1812 年关于超几何级数的论文中（{[}9{]}，第 III 卷，第 139 页），Gauss 给出了收敛性讨论的第一个典范，用他自己的话说，这讨论“以完全的严格性进行，为的是让那些偏爱古代几何学家的严格方法的人满意”。诚然，这一在论文中居于次要地位的讨论并没有追溯到级数理论的第一原理；Cauchy 在其 1821 年的《分析教程》{[}10{]}中，第一个从 Cauchy 判别准则（有清楚的陈述，并被当作自明的）出发，以在每一点上都正确的方式确立了这些原理。由于他在数的定义上采取 Barrow 与 Newton 的观点，可以说，对 Cauchy 而言实数是由量的公理与 Cauchy 判别准则来定义的；而事实上这就足以定义它们（见第五章 § 2）。

(*) 设给定 \(x_0, y_0\) 且 \(> 0\) ，令 \(x_{n+1} = (x_n + y_n)/2\) 以及

\[
y _ {n + 1} = \sqrt {x _ {n} y _ {n}};
\]

与此同时，实数理论的另一个重要方面也得到了最终的澄清。如前所述，两条连续曲线要互相穿过而不相交，在几何上一直被视为显然——这一原理（经适当精确化后）同样等价于实直线（作为一致空间）的完备性。这一原理是 Gauss 于 1799 年给出的 d'Alembert 定理“严格”证明的基础，该定理断言每个实系数多项式都有一个实的或复的根（{[}9{]}，第 III 卷，第 1 页）；Gauss 于 1815 年对同一定理给出的证明（{[}9{]}，第 III 卷，第 31 页）与 Lagrange 较早的一次尝试一样，依赖的是类似而更简单的原理：多项式不可能变号而不取零；这是 § 6 定理 2 的一个特例，我们已经看到 Stevin 使用过它。1817 年，Bolzano 基于这后一原理给出了一个完全的证明，他把这一原理作为关于单实变量的连续实值函数的类似定理的特例而得到{[}11{]}。他清楚地陈述了“Cauchy 判别准则”（并且在 Cauchy 之前），并试图论证它；在缺乏实数的任何算术定义的情况下，他的论证过去只能、也只能是一个恶性循环；但是，一旦接受了这一点，他的工作是完全正确的而且是十分出色的，因为其中不仅包含连续函数的现代定义（在此首次给出）以及多项式连续性的证明，而且还包含任意有界实数集下确界存在的证明（他说的不是集合，而是实数的性质，这二者是一回事）。Cauchy 在其《分析教程》{[}10{]}中也定义了一个或多个实变量的连续函数，并用与 Simon Stevin 相同的推理证明了单变量连续函数不可能变号而不取零。一旦定义了连续性，只要援引 Cauchy 判别准则（或者像 Cauchy 在此处那样，采用等价的“区间套”原理；十进分数的收敛当然只是这一原理的一个特例），这一推理就成为正确的了。

到达这一步之后，留给数学家们的工作只是通过纠正各种错误、填补各种空缺，把所得结果加以精确化并予以推广。例如，Cauchy 曾一度相信，各项都是单变量连续函数的收敛级数，其和是一个连续函数。Abel 在其关于级数的重要工作中（{[}12{]}，第 1 卷，第 219 页；另参看第 2 卷，第 257 页及以下各处）对这一点的纠正，最终导致 Weierstrass 阐明了一致收敛的概念（在其讲义中，这些讲义一直未出版，但影响巨大；见第十章的历史注记）。又如，Cauchy 在其多项式根存在性的一个证明中，未经充分论证就假定了连续函数最小值的存在；而又是 Weierstrass 通过（在其讲义中）对定义在有界闭区间上的实变量函数证明 § 6 的定理 1，阐明了这一类问题。继他对把这一定理不当应用于函数集（“Dirichlet 原理”是最著名的例子）的批评之后，开始了那样一场思想运动，它如我们在第一章的历史注记中所见，导致了紧空间的一般定义以及我们所给出的该定理的现代陈述。

与此同时，Weierstrass 在其讲义中已经察觉到使实数观念完全独立于量理论在逻辑上的重要性；后者实际上等价于对直线上的点（从而对实数集）的公理化定义，外加这样一个集合存在的假定；这种方法虽然本质上是正确的，但显然更好的做法是只从有理数出发，并通过完备化由它们构造出实数（*）。Weierstrass、Dedekind、Méray 与 Cantor 以各不相同的方法、彼此独立地做到了这一点；Dedekind {[}13{]} 提出的“分割”方法与 Eudoxus 的定义十分接近，而所提出的其他方法则都接近于本丛书中所阐述的方法。与此同时，Cantor 开始发展实数集的理论——这一思想最初由 Dedekind 孕育（见第一章的参考文献）——并由此得到了关于实直线的拓扑、其开集与闭集的结构、导集与全不连通完备集等概念的主要初等结果；他还得到了 § 8 关于连续统的势的定理 1，并由它推出：连续统是不可数的，超越数所成之集具有连续统的势，而且（在当时是一个悖论性的结果）平面（或空间）的点集与直线的点集具有相同的势。

随着 Cantor 的工作，本章所研究的问题实际上具备了它们的最终形式。关于 Cantor 工作之直接影响的说明，请参看第一章的历史注记；这里只简要指出它所延伸的方向。除了导致一般拓扑方面的工作（见第一章）以及积分学方面的应用——这两方面本丛书将在别处详尽论述——之外，Cantor 的工作还引发了对直线上点集以及实变量的实值函数的结构与分类的研究。这些研究起源于 Borel {[}14{]} 的工作，后者主要面向测度论，但除其他成果外，它引出了“Borel 集”的定义，即属于 R 的满足下列条件的最小子集族的那些集：该族包含所有区间，且关于并与可数交以及取余运算 C 是封闭的（参看第九章，§ 6，第 3 目）。这些集与所谓“Borel 函数”或“Baire 函数”密切相关，后者是指可以由连续函数通过取序列极限的运算、“超限地”反复进行而得到的函数；Baire 在一项重要的工作中定义了它们，在这项工作中他完全抛弃了测度的观点，转而对这些问题的定性与“拓扑”的方面作系统的考察{[}15{]}；在这项工作中他定义并研究了半连续函数（并且是第一个这样做的人），而为了刻画作为连续函数之极限的函数，他引入了瘦集（Baire 术语中“第一范畴”的集）这一重要概念，我们将在第九章研究它。至于继 Baire 之后的大量工作——它们主要是俄国学派尤其是波兰学派的产物——我们在这里只能提请注意它们的存在（例如可参看{[}16{]}以及 Fundamentae Mathematicae 杂志）。

